\chapter{Panel Econometrics}
\label{ch:panel}

\begin{projectconnection}[Project Connection --- Phase 4A Panel Extension]
In Phase 4A you move from firm-level aggregates to asset-class-level risk outputs and test within-class vs.\ cross-class prediction using panel data.
This chapter teaches the econometric methods you need: fixed effects, clustered standard errors, and Fama--MacBeth regressions.
The panel structure is also how you increase your effective sample size---a critical concern when firm-level VaR gives only ${\sim}1{,}250$ daily observations over five years.
Stacking four asset classes doubles your observations; stacking seven (the He--Kelly--Manela setup) nearly sextuples them.
\end{projectconnection}

% ══════════════════════════════════════════════════════════════
\section{Why Panel Data}
\label{sec:panel-why}
% ══════════════════════════════════════════════════════════════

\begin{definition}[Panel Data]
A \textbf{panel} (or \textbf{longitudinal}) data set tracks multiple entities $i = 1, \dots, N$ over multiple time periods $t = 1, \dots, T$.
Each observation is indexed by the pair $(i, t)$, giving a total of up to $N \times T$ rows.
When every entity is observed at every date, the panel is \emph{balanced}; when some $(i,t)$ pairs are missing, it is \emph{unbalanced}.
\end{definition}

Three data structures appear throughout empirical finance:
\begin{enumerate}[leftmargin=*]
  \item \textbf{Cross-section.} Many entities, one date. Example: returns of 500 stocks on a single day.
  \item \textbf{Time series.} One entity, many dates. Example: daily VIX over five years ($T \approx 1{,}250$).
  \item \textbf{Panel.} Many entities, many dates. Example: daily VaR utilisation for each of $N = 7$ asset classes over $T = 1{,}250$ trading days.
\end{enumerate}

Panels give you three advantages a pure time series cannot:
\begin{itemize}[leftmargin=*]
  \item \textbf{More observations.} Stacking $N$ entities multiplies your sample size. With $N = 7$ asset classes and $T = 1{,}250$ days you have up to $8{,}750$ rows instead of $1{,}250$.
  \item \textbf{Control for unobserved heterogeneity.} Each entity may have a different average level of the dependent variable for reasons you cannot measure. Fixed effects absorb these permanent differences (Section~\ref{sec:panel-fe}).
  \item \textbf{Cross-sectional variation in exposure.} At each date, some asset classes have high VaR utilisation and others have low. This within-date variation identifies the coefficient on VaR utilisation even if aggregate risk is trending over time.
\end{itemize}

\begin{prereq}[Your Panel]
In this project your panel has:
\begin{itemize}[leftmargin=*]
  \item \textbf{Entities} ($i$): asset classes---equities, rates, credit, FX, commodities, and potentially finer subdivisions.
  \item \textbf{Time} ($t$): trading days, typically $T \approx 1{,}000$--$1{,}500$ depending on history available in SecDB.
  \item \textbf{Features} ($\bx_{it}$): asset-class-specific VaR utilisation, factor concentration ($\HHI$), scenario P\&L rank, component VaR share.
  \item \textbf{Targets} ($y_{it}$): next-period return, realised volatility, drawdown indicator, or VIX innovation.
\end{itemize}
The key question is whether a risk signal that predicts returns in \emph{one} asset class also predicts returns in \emph{other} asset classes---the ``single pricing kernel'' test of He--Kelly--Manela (2017).
\end{prereq}

% ══════════════════════════════════════════════════════════════
\section{Fixed Effects}
\label{sec:panel-fe}
% ══════════════════════════════════════════════════════════════

\begin{prereq}[Dummy Variables and Degrees of Freedom]
A \textbf{dummy variable} (indicator variable) equals 1 for one category and 0 otherwise.
Including $N$ dummies for $N$ groups, plus an intercept, creates perfect collinearity; drop one dummy or suppress the intercept.
Each dummy costs one degree of freedom, so adding 7 asset-class dummies to a regression with $8{,}750$ observations reduces $df$ by 6---negligible.
With thousands of firm-level dummies, the cost is material.
\end{prereq}

\begin{definition}[Fixed Effects Model]
The \textbf{entity fixed-effects} (FE) model is:
\begin{equation}
  y_{it} = \alpha_i + \bx_{it}'\bbeta + \varepsilon_{it}
  \label{eq:fe-model}
\end{equation}
where:
\begin{itemize}[leftmargin=*]
  \item $y_{it}$ --- outcome for entity $i$ at time $t$ (e.g., next-day return for asset class $i$),
  \item $\alpha_i$ --- entity-specific intercept, absorbing all \emph{time-invariant} differences across entities,
  \item $\bx_{it}$ --- $K \times 1$ vector of time-varying regressors (e.g., VaR utilisation, factor $\HHI$),
  \item $\bbeta$ --- $K \times 1$ coefficient vector, common across entities,
  \item $\varepsilon_{it}$ --- idiosyncratic error, assumed $\E[\varepsilon_{it} \mid \alpha_i, \bx_{it}] = 0$.
\end{itemize}
\end{definition}

The intercept $\alpha_i$ captures everything that differs across entities but does not change over time: equities have structurally higher average returns than FX carry; credit VaR operates at different levels than rates VaR.
Fixed effects absorb these level differences so that $\bbeta$ is identified purely from \emph{within-entity} variation over time.

\subsection{The Within Estimator}

\begin{keyidea}[Demeaning Removes Entity Effects]
To estimate $\bbeta$ in the FE model, subtract each entity's time-series mean from both sides:
\begin{equation}
  y_{it} - \bar{y}_i = (\bx_{it} - \bar{\bx}_i)'\bbeta + (\varepsilon_{it} - \bar{\varepsilon}_i)
  \label{eq:within-estimator}
\end{equation}
where $\bar{y}_i = \frac{1}{T}\sum_{t=1}^{T} y_{it}$ and similarly for $\bar{\bx}_i$ and $\bar{\varepsilon}_i$.
The entity effects $\alpha_i$ cancel.
OLS on the demeaned data yields the \textbf{within estimator}, also called the \textbf{LSDV} (Least Squares Dummy Variable) estimator.
\end{keyidea}

You can equivalently include $N$ dummy variables---one per entity---and run OLS on the full model.
The coefficient $\hat{\bbeta}$ is identical; demeaning is just computationally cheaper when $N$ is large (Wooldridge, 2010, Chapter~10).

\begin{prereq}[Between vs.\ Within Variation]
Panel data contains two sources of variation:
\begin{itemize}[leftmargin=*]
  \item \textbf{Between variation}: differences \emph{across} entities in their time-series averages. Example: equities have higher average VaR utilisation than FX.
  \item \textbf{Within variation}: fluctuations \emph{over time within} a single entity. Example: equity VaR utilisation rises from 60\% to 85\% during a stress event, then falls back.
\end{itemize}
The FE estimator discards between variation entirely---it uses only within variation.
This is a feature, not a bug: between variation is contaminated by omitted time-invariant confounders.
The cost is that you cannot estimate the effect of any time-invariant regressor (e.g., ``is this asset class a derivative?'').
If you need those effects, consider Mundlak's (1978) compromise: include entity means $\bar{\bx}_i$ as additional regressors in a RE model, which is algebraically equivalent to FE for the time-varying coefficients while also estimating between-entity effects.
\end{prereq}

\begin{workedexample}{Fixed Effects with 4 Asset Classes}
Suppose you have $N = 4$ asset classes (equities, rates, credit, FX) observed over $T = 252$ trading days, giving $4 \times 252 = 1{,}008$ panel observations.
You regress next-day return $y_{it}$ on VaR utilisation $x_{it}$ with entity fixed effects:
\[
  y_{it} = \alpha_i + \beta \cdot x_{it} + \varepsilon_{it}
\]
\textbf{Step 1: Compute entity means.}
For each asset class $i$, compute $\bar{y}_i$ and $\bar{x}_i$ across its 252 days.

\textbf{Step 2: Demean.}
Create $\tilde{y}_{it} = y_{it} - \bar{y}_i$ and $\tilde{x}_{it} = x_{it} - \bar{x}_i$.

\textbf{Step 3: OLS on demeaned data.}
\[
  \hat{\beta}_{\text{FE}} = \frac{\sum_{i=1}^{4}\sum_{t=1}^{252} \tilde{x}_{it}\,\tilde{y}_{it}}{\sum_{i=1}^{4}\sum_{t=1}^{252} \tilde{x}_{it}^2}
\]
If $\hat{\beta}_{\text{FE}} = -0.03$ with $\text{SE} = 0.01$, the interpretation is: within a given asset class, a one-standard-deviation increase in VaR utilisation predicts a 3\,bp decrease in next-day return.
The negative sign is consistent with the Coval--Stafford (2007) forced-selling channel (see Chapter~\ref{ch:intermediary}).

\textbf{Degrees of freedom:} $N \times T - N - K = 1{,}008 - 4 - 1 = 1{,}003$.
The 4 dummies consume only 4 degrees of freedom.
\end{workedexample}

\subsection{Time Fixed Effects}

You can also include \textbf{time fixed effects} $\delta_t$:
\begin{equation}
  y_{it} = \alpha_i + \delta_t + \bx_{it}'\bbeta + \varepsilon_{it}
  \label{eq:twoway-fe}
\end{equation}
Time effects absorb market-wide shocks common to all entities on a given day---aggregate volatility spikes, macro announcements, or Fed meetings.
This is a \textbf{two-way fixed effects} model.
After adding time effects, $\bbeta$ is identified only from variation that differs \emph{across entities within the same date}, net of the entity average.

\begin{prereq}[When to Use Time Fixed Effects]
Include time fixed effects when you want to control for aggregate shocks that hit all asset classes simultaneously.
This is appropriate when you care about the \emph{relative} effect of VaR utilisation across asset classes, not whether VaR utilisation predicts returns in general.
If your goal is to predict absolute returns, omit time effects---they absorb the common signal.
\end{prereq}

% ══════════════════════════════════════════════════════════════
\section{Random Effects and the Hausman Test}
\label{sec:panel-re}
% ══════════════════════════════════════════════════════════════

The fixed-effects model treats $\alpha_i$ as parameters to estimate.
The \textbf{random effects} (RE) model treats them as draws from a distribution:
\begin{equation}
  y_{it} = \mu + \bx_{it}'\bbeta + \alpha_i + \varepsilon_{it}, \qquad
  \alpha_i \sim (0, \sigma_\alpha^2), \quad
  \varepsilon_{it} \sim (0, \sigma_\varepsilon^2)
  \label{eq:re-model}
\end{equation}
The critical assumption is $\E[\alpha_i \mid \bx_{it}] = 0$: the entity effect is \emph{uncorrelated} with the regressors.

If the RE assumption holds, the RE estimator is more efficient than FE because it uses both within-entity and between-entity variation.
If the assumption fails---$\alpha_i$ correlates with $\bx_{it}$---the RE estimator is inconsistent, and FE is the safe choice.

\begin{intuition}[Fixed Effects vs.\ Random Effects]
Fixed effects says: ``I do not care \emph{why} asset classes differ in their average return levels.
Maybe equities have a higher equity premium, maybe credit has illiquidity compensation---I just want to control for it.''
Random effects says: ``The differences across asset classes are random draws from a common distribution, and they are unrelated to VaR utilisation.''

In finance, the RE assumption almost always fails.
Asset classes with structurally higher returns (equities) also tend to have structurally different risk profiles (higher VaR utilisation limits).
That correlation between $\alpha_i$ and $\bx_{it}$ biases the RE estimator.
Default to fixed effects.
\end{intuition}

\subsection{The Hausman Test}

The \textbf{Hausman test} (1978) formalises the FE-vs-RE choice:
\begin{equation}
  H = (\hat{\bbeta}_{\text{FE}} - \hat{\bbeta}_{\text{RE}})' \bigl[\widehat{\operatorname{Var}}(\hat{\bbeta}_{\text{FE}}) - \widehat{\operatorname{Var}}(\hat{\bbeta}_{\text{RE}})\bigr]^{-1} (\hat{\bbeta}_{\text{FE}} - \hat{\bbeta}_{\text{RE}}) \;\sim\; \chi^2_K
  \label{eq:hausman}
\end{equation}
where $K$ is the number of regressors.
Under $H_0$: RE is consistent (both estimators converge to the same $\bbeta$).
Under $H_1$: RE is inconsistent (FE and RE diverge).
If you reject $H_0$, use FE.

In practice, for finance panels, the Hausman test almost always rejects.
Use fixed effects as your default.

% ══════════════════════════════════════════════════════════════
\section{Clustered Standard Errors}
\label{sec:panel-cluster}
% ══════════════════════════════════════════════════════════════

Fixed effects handle the \emph{point estimate} correctly, but the \emph{standard errors} require separate treatment.
The problem: panel residuals are not independent.
Returns of the same asset class are correlated over time (autocorrelation), and returns of different asset classes on the same date are correlated in the cross section (common shocks).
Ignoring this correlation produces standard errors that are too small, $t$-statistics that are too large, and false discoveries.

\begin{keyresult}[Petersen (2009) --- Standard Errors in Finance Panels]
Petersen (2009, \textit{Review of Financial Studies}) uses Monte Carlo simulations to show that OLS standard errors are severely downward biased in finance panels.
Specific findings:
\begin{itemize}[leftmargin=*]
  \item With a firm effect ($\rho_X = \rho_\varepsilon = 0.50$), OLS reports $\widehat{\text{SE}} = 0.0283$ while the true standard error is $0.0506$---understating by 44\%.
  \item The bias grows with panel length: from ${\sim}30\%$ understatement with 5 years per entity to ${\sim}73\%$ with 50 years.
  \item Using unclustered OLS standard errors, 15.3\% of $t$-statistics are significant at the 1\% level---over 15 times the nominal rate.
  \item Only clustering by the correct dimension (firm, time, or both) produces unbiased standard errors.
\end{itemize}
\end{keyresult}

\begin{warning}[Never Report Unclustered Standard Errors from a Finance Panel]
Unclustered standard errors in finance panels produce $t$-statistics that are inflated by a factor that grows with the panel's length and the strength of within-cluster correlation.
In Petersen's simulations, a nominal 1\% test rejects 15.3\% of the time.
A ``significant'' result with unclustered SEs is not evidence of anything.
Always cluster---the question is only \emph{by what dimension}.
\end{warning}

\subsection{One-Way Clustering}

\begin{definition}[Clustered Standard Errors]
Let the panel regression be $y_{it} = \bx_{it}'\bbeta + \varepsilon_{it}$.
The \textbf{cluster-robust} variance estimator, clustering by entity $i$, is:
\begin{equation}
  \widehat{\operatorname{Var}}(\hat{\bbeta}) = (\bX'\bX)^{-1} \left(\sum_{i=1}^{N} \bX_i' \hat{\bm{\varepsilon}}_i \hat{\bm{\varepsilon}}_i' \bX_i \right) (\bX'\bX)^{-1}
  \label{eq:cluster-var}
\end{equation}
where $\bX_i$ is the $T \times K$ submatrix of regressors for entity $i$ and $\hat{\bm{\varepsilon}}_i$ is the $T \times 1$ vector of residuals for entity $i$.
\end{definition}

\textbf{Cluster by entity} when residuals are correlated within the same entity over time (autocorrelation).
This is the standard in corporate finance panels where firms have persistent unobservables.

\textbf{Cluster by time} when residuals are correlated across entities within the same period (cross-sectional dependence).
This is common in asset pricing when a market shock hits all assets on the same date.

\subsection{Two-Way (Double) Clustering}

Finance panels typically exhibit \emph{both} entity-level and time-level correlation.
The two-way cluster-robust variance estimator (Cameron, Gelbach, and Miller, 2011) is:
\begin{equation}
  \widehat{\operatorname{Var}}_{\text{double}} = \widehat{\operatorname{Var}}_{\text{entity}} + \widehat{\operatorname{Var}}_{\text{time}} - \widehat{\operatorname{Var}}_{\text{entity} \times \text{time}}
  \label{eq:double-cluster}
\end{equation}
where $\widehat{\operatorname{Var}}_{\text{entity} \times \text{time}}$ clusters by the intersection (each unique $(i,t)$ pair---equivalent to heteroskedasticity-robust SEs).
Petersen (2009) shows this is the only approach that produces correct standard errors when both dimensions of correlation are present.

\begin{workedexample}{Unclustered vs.\ Clustered Standard Errors}
You run a panel regression of next-day return on VaR utilisation with entity fixed effects, using $N = 7$ asset classes and $T = 1{,}000$ days ($7{,}000$ observations).
The coefficient estimate is $\hat{\beta} = -0.025$.

\medskip
\begin{center}
\begin{tabular}{lccc}
  \toprule
  \textbf{SE Method} & \textbf{SE($\hat{\beta}$)} & \textbf{$t$-statistic} & \textbf{$p$-value} \\
  \midrule
  OLS (no clustering)          & 0.004 & $-6.25$ & $< 0.001$ \\
  Clustered by entity          & 0.008 & $-3.13$ & $0.002$ \\
  Clustered by time            & 0.007 & $-3.57$ & $< 0.001$ \\
  Double-clustered             & 0.010 & $-2.50$ & $0.013$ \\
  \bottomrule
\end{tabular}
\end{center}
\medskip

The unclustered $t$-statistic ($-6.25$) overstates significance by a factor of 2.5 relative to the double-clustered $t$-statistic ($-2.50$).
In this example all methods still reject at 5\%, but in a marginal case the unclustered SE would produce a false discovery.
With $N = 7$ clusters on the entity dimension, entity-clustered SEs may themselves be unreliable---the asymptotic approximation requires $N$ large.
When $N$ is small, report double-clustered SEs and note the small $N$ caveat, or use the wild cluster bootstrap (Cameron, Gelbach, and Miller, 2008).
\end{workedexample}

\begin{prereq}[How Many Clusters Is Enough?]
Clustered standard errors rely on an asymptotic argument in the number of clusters.
With $N = 7$ asset classes, you have only 7 entity-level clusters---below the common rule-of-thumb of ${\sim}50$ for reliable asymptotics (Cameron and Miller, 2015).
Solutions: (a) use finer entities (sub-desks or individual products rather than broad asset classes); (b) use the wild cluster bootstrap; (c) cluster by time (where $T \approx 1{,}000$ gives many clusters) and note the limitation.
\end{prereq}

% ══════════════════════════════════════════════════════════════
\section{Fama--MacBeth Regressions}
\label{sec:panel-fm}
% ══════════════════════════════════════════════════════════════

Fama and MacBeth (1973) introduced an alternative to pooled panel regression that has become the workhorse of cross-sectional asset pricing.

\begin{keyidea}[Fama--MacBeth Two-Pass Procedure]
\textbf{Pass 1 (time-series).} For each asset $i = 1, \dots, N$, regress returns on factor realisations to estimate factor loadings:
\begin{equation}
  R_{it} = \alpha_i + \beta_{i,1} F_{1t} + \cdots + \beta_{i,K} F_{Kt} + \varepsilon_{it}, \qquad t = 1, \dots, T
  \label{eq:fm-pass1}
\end{equation}
This gives $\hat{\beta}_{i,k}$ for each asset and factor.

\textbf{Pass 2 (cross-sectional).} At each date $t$, regress returns on the estimated betas:
\begin{equation}
  R_{it} = \gamma_{0t} + \gamma_{1t} \hat{\beta}_{i,1} + \cdots + \gamma_{Kt} \hat{\beta}_{i,K} + u_{it}, \qquad i = 1, \dots, N
  \label{eq:fm-pass2}
\end{equation}
This gives a time series of coefficient estimates $\{\hat{\gamma}_{kt}\}_{t=1}^T$ for each factor $k$.

\textbf{Estimate and inference.} The risk premium for factor $k$ is the time-series average:
\begin{equation}
  \hat{\lambda}_k = \frac{1}{T}\sum_{t=1}^T \hat{\gamma}_{kt}
  \label{eq:fm-lambda}
\end{equation}
The standard error uses the time-series variation of the cross-sectional estimates:
\begin{equation}
  \widehat{\text{SE}}(\hat{\lambda}_k) = \frac{1}{\sqrt{T}} \sqrt{\frac{1}{T-1}\sum_{t=1}^T (\hat{\gamma}_{kt} - \hat{\lambda}_k)^2}
  \label{eq:fm-se}
\end{equation}
\end{keyidea}

The key insight is that by running separate cross-sectional regressions at each date and then averaging, Fama--MacBeth \emph{automatically accounts for cross-sectional correlation}.
If all assets move together on a given date, the cross-sectional coefficient for that date reflects it, and the time-series averaging smooths it out.

However, Fama--MacBeth does \emph{not} correct for time-series autocorrelation.
If $\hat{\gamma}_{kt}$ is persistent over time, the standard error in Equation~\eqref{eq:fm-se} is too small.
The fix: apply Newey--West (1987) correction to the time series of $\hat{\gamma}_{kt}$ estimates.

\begin{prereq}[Fama--MacBeth vs.\ Clustered Panel Regression]
Petersen (2009) proves:
\begin{itemize}[leftmargin=*]
  \item Fama--MacBeth SEs $\approx$ clustering by time. Both handle cross-sectional correlation; neither handles time-series correlation unless adjusted.
  \item Clustering by entity handles time-series correlation but not cross-sectional correlation.
  \item Double-clustering handles both.
\end{itemize}
Use Fama--MacBeth for \textbf{pricing tests} (``is this factor priced in the cross section?'').
Use panel FE with double-clustered SEs for \textbf{prediction} (``does VaR utilisation predict returns, controlling for entity effects?'').
\end{prereq}

\subsection{Fama--MacBeth in Practice: He--Kelly--Manela}

He, Kelly, and Manela (2017, \textit{JFE}) use Fama--MacBeth regressions to test whether intermediary capital risk is priced across asset classes.
Their setup:

\begin{itemize}[leftmargin=*]
  \item \textbf{Factor:} Shocks to the equity capital ratio of New York Fed primary dealers.
  \item \textbf{Test assets:} Portfolios spanning seven asset classes---equities, US government bonds, corporate bonds, sovereign bonds, options, CDS, commodities, and FX.
  \item \textbf{Pass 1:} Estimate each portfolio's beta to the intermediary capital factor using time-series regression.
  \item \textbf{Pass 2:} Cross-sectional regression of average excess returns on estimated betas.
\end{itemize}

\begin{keyresult}[He--Kelly--Manela (2017) --- Single Pricing Kernel]
A single intermediary capital risk factor prices the cross section of returns across seven asset classes with a cross-sectional $R^2$ of 77\%.
The estimated capital risk premium is approximately 9\% per year, with a range of 7--11\% for five of the seven asset classes.
The price of risk is consistently positive and of similar magnitude when estimated separately for individual asset classes---evidence for a single stochastic discount factor driven by dealer balance-sheet constraints.
\end{keyresult}

This is the test you replicate in Phase 4A: if SecDB risk outputs capture the same dealer-constraint channel, your risk signals should price assets across classes in a similar Fama--MacBeth framework.

\begin{workedexample}{Fama--MacBeth with 3 Asset Classes}
Suppose you have $N = 20$ portfolios across 3 asset classes and $T = 60$ monthly observations.
You want to test whether exposure to a VaR-utilisation factor earns a risk premium.

\textbf{Pass 1.} For each portfolio $i$, run a time-series regression:
\[
  R_{it} = \alpha_i + \beta_i \cdot \Delta\text{VaRUtil}_t + \varepsilon_{it}, \qquad t = 1, \dots, 60
\]
This gives 20 estimated betas $\hat{\beta}_1, \dots, \hat{\beta}_{20}$.

\textbf{Pass 2.} At each month $t$, regress the cross section of 20 portfolio returns on their betas:
\[
  R_{it} = \gamma_{0t} + \gamma_{1t}\,\hat{\beta}_i + u_{it}, \qquad i = 1, \dots, 20
\]
This gives 60 estimates of $\hat{\gamma}_{1t}$---the risk premium for VaR-utilisation exposure in month $t$.

\textbf{Estimate.} Average: $\hat{\lambda} = \frac{1}{60}\sum_{t=1}^{60}\hat{\gamma}_{1t} = 0.42\%$ per month.

\textbf{Standard error.}
$\widehat{\text{SE}}(\hat{\lambda}) = \frac{1}{\sqrt{60}}\sqrt{\frac{1}{59}\sum_t (\hat{\gamma}_{1t} - 0.42\%)^2} = 0.18\%$.

\textbf{$t$-statistic.} $t = 0.42 / 0.18 = 2.33$ ($p = 0.023$).

Interpretation: portfolios with higher exposure to VaR-utilisation shocks earn 42\,bp/month (5.0\% annualised) more than low-exposure portfolios, significant at 5\%.
\end{workedexample}

% ══════════════════════════════════════════════════════════════
\section{Panel ML vs.\ Panel Regression}
\label{sec:panel-ml}
% ══════════════════════════════════════════════════════════════

Panel regression and ML are complements, not substitutes.
They answer different questions:

\begin{center}
\begin{tabular}{lll}
  \toprule
  & \textbf{Panel FE Regression} & \textbf{Panel ML (LightGBM)} \\
  \midrule
  \textbf{Goal} & Causal interpretation, pricing tests & Out-of-sample prediction \\
  \textbf{Inference} & Clustered SEs, $t$-statistics & Purged CV, information coefficient \\
  \textbf{Flexibility} & Linear in features & Nonlinear, interactions \\
  \textbf{Entity effects} & Absorbed via demeaning & Dummy features for asset class \\
  \textbf{Interpretability} & Coefficient table & $\SHAP$ values \\
  \textbf{Weakness} & Misses nonlinearities & No valid $p$-values \\
  \bottomrule
\end{tabular}
\end{center}

\subsection{Using Both in Your Project}

\begin{keyidea}[Two-Track Panel Analysis]
Your project uses panel data in two complementary ways:
\begin{enumerate}[leftmargin=*]
  \item \textbf{Panel FE regression} for the He--Kelly--Manela pricing test.
    Run Fama--MacBeth or panel FE with double-clustered SEs.
    Report $\hat{\lambda}$ (price of risk), $t$-statistics, and cross-sectional $R^2$.
    This tells you whether SecDB risk signals are \emph{priced}---whether bearing exposure to VaR-utilisation shocks earns a premium.

  \item \textbf{LightGBM with asset-class indicators} for prediction.
    Include one-hot encoded asset-class dummies as features alongside risk signals.
    The tree can learn asset-class-specific nonlinear relationships (e.g., VaR utilisation matters more for credit than for FX).
    Evaluate with purged cross-validation (Chapter~\ref{ch:validation}) and $\IC$ (Chapter~\ref{ch:backtesting}).
\end{enumerate}
The regression answers ``is the signal priced?'' with proper statistical inference.
The ML model answers ``can the signal predict?'' with proper out-of-sample evaluation.
Present both in your deliverable.
\end{keyidea}

\subsection{Entity Features in Tree Models}

When you include asset-class dummies in LightGBM, the tree can split on them.
A split on ``is\_credit = 1'' followed by a split on VaR utilisation is equivalent to estimating a different VaR-utilisation slope for credit vs.\ other classes---an interaction effect that panel FE with a single $\bbeta$ cannot capture without explicit interaction terms.

However, with only $N = 7$ asset classes, the dummies are low-cardinality.
LightGBM handles this natively.
Do \emph{not} target-encode the asset-class variable---with 7 categories and time-series data, target encoding will leak future information unless you are extremely careful with temporal ordering.

\begin{prereq}[Panel ML Validation Pitfall]
When cross-validating a panel model, you must respect the time dimension.
A random $K$-fold split will place observations from the same date into both train and test folds, leaking cross-sectional correlation.
Use the purged $K$-fold CV from Chapter~\ref{ch:validation}: split by \emph{time}, not by row.
For panel data, this means entire dates go into train or test, never split across folds.
\end{prereq}

% ══════════════════════════════════════════════════════════════
\section{Choosing the Right Panel Method}
\label{sec:panel-choice}
% ══════════════════════════════════════════════════════════════

With multiple methods available, the decision tree is:

\begin{enumerate}[leftmargin=*]
  \item \textbf{Is your question ``is this factor priced?'' or ``can I predict returns?''}
    \begin{itemize}[leftmargin=*]
      \item Pricing test $\Rightarrow$ Fama--MacBeth (Section~\ref{sec:panel-fm}). Report $\hat{\lambda}$, $t$-stat, cross-sectional $R^2$.
      \item Prediction $\Rightarrow$ Panel FE regression or panel ML (Sections~\ref{sec:panel-fe},~\ref{sec:panel-ml}).
    \end{itemize}
  \item \textbf{Do you need interpretable coefficients with valid inference?}
    \begin{itemize}[leftmargin=*]
      \item Yes $\Rightarrow$ Panel FE with double-clustered SEs. Report $\hat{\bbeta}$, $t$-stats.
      \item No, you want the best forecast $\Rightarrow$ LightGBM with asset-class dummies and purged CV.
    \end{itemize}
  \item \textbf{How many entity-level clusters do you have?}
    \begin{itemize}[leftmargin=*]
      \item $N \geq 50$ $\Rightarrow$ entity clustering is reliable.
      \item $N < 50$ $\Rightarrow$ cluster by time ($T$ is large), or use wild cluster bootstrap, and note the caveat.
    \end{itemize}
  \item \textbf{Present both.} In your deliverable, the panel FE regression gives the desk a causal story (``VaR utilisation predicts returns because forced selling depletes liquidity''), while the ML model gives a tradeable signal.
\end{enumerate}

% ══════════════════════════════════════════════════════════════
\section{Practical Considerations}
\label{sec:panel-practical}
% ══════════════════════════════════════════════════════════════

\subsection{Unbalanced Panels}

Your panel will likely be unbalanced: CDS data starts later than equity data, some products were introduced or discontinued, and SecDB history depth varies by desk.
Fixed effects handle unbalanced panels naturally---the within estimator only requires variation over whatever time periods exist for each entity.
Document which entities have shorter histories and check that results are not driven solely by the longest series.

\subsection{Stationarity}

Panel regressions assume stationarity within each entity's time series.
If VaR utilisation has a unit root (it drifts without mean reversion), the regression is spurious.
Pre-test each feature: compute the augmented Dickey--Fuller statistic within each entity.
If a feature is non-stationary, take first differences or use the $z$-score (deviation from rolling mean divided by rolling standard deviation), which you already construct in your feature engineering pipeline (Chapter~\ref{ch:features}).

\subsection{Survivorship Bias}

If your panel includes only asset classes that currently exist in SecDB, you miss any that were discontinued.
This is less of a concern for broad asset classes (equities and rates are permanent) than for individual products.
Note the limitation; do not overstate generalisability.

\subsection{Implementation in Python}

\begin{prereq}[Python Packages for Panel Econometrics]
The \texttt{linearmodels} package (by Kevin Sheppard) implements panel FE, RE, Fama--MacBeth, and clustered SEs directly:
\begin{itemize}[leftmargin=*]
  \item \texttt{PanelOLS} --- entity/time fixed effects with clustered SEs.
  \item \texttt{FamaMacBeth} --- two-pass Fama--MacBeth with Newey--West adjustment.
  \item \texttt{RandomEffects} --- RE estimator with Hausman test via \texttt{compare()}.
\end{itemize}
Set your DataFrame to a \texttt{MultiIndex} with levels \texttt{(entity, time)} before fitting.
Specify clustering with the \texttt{cluster\_entity} and \texttt{cluster\_time} arguments.
Always check that the index is sorted by time within each entity to avoid silent misalignment.

For the ML track, \texttt{LightGBM} does not know about panel structure---it treats each row independently.
You enforce temporal ordering through your validation strategy (purged $K$-fold from Chapter~\ref{ch:validation}), not through the model itself.
\end{prereq}

\subsection{Nickell Bias}

When $T$ is small relative to the persistence of the dependent variable, the within estimator is biased.
This is the \textbf{Nickell (1981) bias}: in a dynamic panel where $y_{it}$ depends on $y_{i,t-1}$, demeaning creates correlation between the demeaned lagged dependent variable and the demeaned error.
The bias is $O(1/T)$, so it vanishes as $T \to \infty$.
With $T \approx 1{,}000$ daily observations, Nickell bias is negligible for your panel.
It becomes a concern only if you aggregate to monthly frequency ($T \approx 60$) and include a lagged dependent variable.

% ══════════════════════════════════════════════════════════════
\section{Summary}
\label{sec:panel-summary}
% ══════════════════════════════════════════════════════════════

\begin{center}
\begin{tabular}{lp{10cm}}
  \toprule
  \textbf{Concept} & \textbf{Key Takeaway} \\
  \midrule
  Panel data & Multiple entities $\times$ multiple time periods; increases sample size and controls for unobserved heterogeneity. \\
  Fixed effects & Absorb entity-level intercepts via demeaning; identify $\bbeta$ from within-entity variation only. \\
  Random effects & More efficient \emph{if} $\alpha_i \perp \bx_{it}$; use Hausman test, but default to FE in finance. \\
  Clustered SEs & OLS SEs are 30--73\% too small in Petersen's (2009) simulations; cluster by entity, time, or both. \\
  Fama--MacBeth & Two-pass: time-series betas, then cross-sectional risk premia; accounts for cross-sectional but not time-series correlation. \\
  He--Kelly--Manela & Single intermediary capital factor prices 7 asset classes with $R^2 = 77\%$ and ${\sim}9\%$ annual premium. \\
  Panel ML & Complement regression with LightGBM using asset-class dummies; validate with time-based purged CV. \\
  \bottomrule
\end{tabular}
\end{center}
